%% 07-the-actuator.tex -- generated by tools/build_from_docx.py; edit freely

\hypertarget{the-actuator}{%
\section{The Actuator}\label{the-actuator}}

The Actuator is the fourth entity of the tuple and the only one whose
output leaves the platform. Monitoring observes, analytics infers and
decision forms intent; all three produce records. The Actuator changes
the network. Everything the framework asserts about containment is
asserted finally here, because an agent contained everywhere else and
unconstrained at the last mile is not contained at all.

Two positions govern the whole chapter. The first is that the Actuator
is a role and not a privilege tier: it executes for normative and
extension agents alike, through the same controller path, and
containment lives in declared control scope, action-space bounds and
mediation rather than in a separate or restricted execution route. A
framework that gave experimental agents a weaker execution path would be
asserting that the path itself carries safety, which would leave the
production path unexamined. What makes a common path acceptable is that
responsibility for an agent's behavior rests with the partner who
contributed it, and that this responsibility is evidenced rather than
assumed: every actuation transaction records the agent instance that
performed it, so each change to the network is attributable to a named
contributor and to the descriptor under which it was admitted. The
platform bounds what an agent may reach; it does not assume the
correctness of the logic inside it.

The second is that an acknowledgement is never proof of change. This is
not a cautious posture but a finding: a transport driver examined in the
consortium's own codebase acknowledges configuration keys it does not
recognize without effecting any physical change, and accepts monitoring
subscriptions without collecting anything. A framework that treated
acknowledgement as success would report healthy service over a network
that had done nothing. Verification against physical state is therefore
structural in this chapter rather than advisory.

\hypertarget{role-and-scope}{%
\subsection{Role and Scope}\label{role-and-scope}}

The Actuator translates intent into enforceable action against domain
controllers, verifies the result against state observed independently of
the controller's answer, and publishes typed execution feedback. It
holds no objectives of its own. It does not choose between candidate
actions, does not defer an action it finds unwise and does not
substitute one action for another; where the intent it receives cannot
be translated within the agent's registered bounds, the transaction
terminates as a rejection with a reason code rather than as an
approximation.

Its core functions are intent-to-action translation, pre-condition state
checks, command staging, guardrail enforcement at the last mile, rate
limiting, idempotency handling, read-back verification, rollback
planning and the publication of execution feedback. Of these, only
translation is contributed logic in the general case. Rate limiting,
retry counting, bound re-validation and feedback emission are held by
the platform runtime, for the same reason they are held by the runtime
elsewhere in the framework: a control enforced by the code it constrains
is not a control.

\hypertarget{controller-bindings}{%
\subsection{Controller Bindings}\label{controller-bindings}}

The Actuator reaches the network through a domain controller and never
through a device directly. In transport the binding is to TeraFlowSDN,
which submits configuration and drives the device. In the radio domain
it is to the near-real-time controller over E2, using the service model
for radio control where near-real-time control is required, and to the
O1 interface for configuration management. At orchestration level the
binding is to the domain management function and the orchestrator's
lifecycle operations.

Every Actuator declares its controller binding and the interface version
of that controller in its interface profile, with a pinned compatibility
range. A binding whose observed version falls outside that range is a
binding fault: actuation is blocked and the fault is raised. It is not
answered with a best-effort call. A best-effort call against an
interface whose semantics have shifted is precisely the condition that
produces an acknowledgement without an effect, and the framework
declines to generate that condition deliberately having found it
occurring by accident.

\hypertarget{capability-declarations-3}{%
\subsection{Capability Declarations}\label{capability-declarations-3}}

An Actuator capability is declared in one of four modes. In full the
capability performs its own translation and execution across its
declared action space. In partial it executes an enumerated subset of
that action space, and an action outside the enumeration is rejected
rather than attempted. In delegated the execution is performed by a
domain controller or another agent, which is named in the descriptor and
which must itself satisfy the requirements of this chapter. In none the
capability does not actuate.

Delegation is the case that most often hides an unexamined path. A
delegate named in the descriptor is admitted on its own terms, and the
requirements of this chapter apply to it as they apply to any Actuator,
so that delegation transfers the work and never the obligations.

\hypertarget{the-four-outcomes}{%
\subsection{The Four Outcomes}\label{the-four-outcomes}}

Every actuation transaction terminates in exactly one of four outcomes,
and the set replaces the success-and-failure binary that most management
interfaces present.

An outcome of acknowledged and verified means the controller
acknowledged the command and an independent read-back confirms that
device state matches the intent. It is the only success state in the
framework.

An outcome of acknowledged but unverified means the controller
acknowledged the command while read-back was absent, inconclusive or
mismatched. It is a failure pending confirmation. It is never reported
as success, never counted toward an objective as achieved, and never
treated as evidence that the network changed.

An outcome of rejected means the controller explicitly refused the
command. It is the cleanest of the failure states: nothing was applied,
the state of the network is known, and retry under the declared policy
is legitimate.

An outcome of unresolved means the command timed out or transport was
lost. The state of the network is unknown and the command may still be
in flight. A timeout is an unresolved transaction and not permission to
retry, because a blind retry from unknown state risks applying a change
twice.

The asymmetry among the four is deliberate. Two of them leave state
known, one leaves it known but wrong, and one leaves it unknown; the
behavior the platform permits in each case follows from which of those
conditions holds, rather than from how many attempts have been made.

\hypertarget{read-back-is-independent-by-construction}{%
\subsubsection{Read-back is independent by construction}\label{read-back-is-independent-by-construction}}

Verification reads network or device state through a path independent of
the acknowledgement: an observation supplied by the Monitoring System,
or a controller state query that reflects applied configuration rather
than request receipt. A controller's own confirmation that it received a
request is not verification of that request, and the distinction is the
whole substance of the outcome set.

Where a controller cannot expose verifiable state for a given action,
the action is declared unverifiable in the action space. Such an action
is admitted with its targets enumerated individually rather than by
class, with a rate ceiling no higher than the one its catalog entry sets
for unverifiable execution, with every execution reported to the domain
management function as it occurs, and under elevated evaluation
scrutiny. Its terminal state is acknowledged but unverified by
construction, and it can never be reported as verified. These
restrictions bound where and how often such an action is taken; they do
not lower the authority of the agent taking it.

\hypertarget{retry-rollback-and-idempotency}{%
\subsection{Retry, Rollback and Idempotency}\label{retry-rollback-and-idempotency}}

Retry safety depends on whether a command may already have applied.
Attempt count alone never establishes that; only declared idempotency
and read-back can. The platform therefore derives retry behavior from
the declaration rather than from operator preference, and enforces it in
the runtime.

From a rejection, retry is permitted under the declared policy, because
state is known and nothing was applied. From an unresolved outcome,
read-back comes first: retry follows only if read-back shows the command
did not apply, and otherwise the command is treated as applied and
proceeds to verification. From an acknowledged but unverified outcome,
retry is permitted only where the command is declared idempotent, and
otherwise the transaction holds and escalates. A non-idempotent command
against unknown state is never retried automatically under any declared
policy.

Rollback is itself an actuation, and performed into unknown state it
compounds damage rather than containing it. Rollback is therefore
permitted only from a state that has been read back. From an unresolved
outcome rollback is prohibited until read-back resolves the actual
state, and where read-back cannot resolve it the Actuator holds and
escalates. Where a declaration is absent or ambiguous, holding and
escalating is the default.

Retry counters, backoff and rate limits are held and applied by the
actuation path rather than by contributed code. An agent that declares
three attempts and whose implementation attempts thirty has the fourth
and subsequent attempts suppressed and a conformance fault recorded
against it.

\hypertarget{last-mile-guardrails}{%
\subsection{Last-Mile Guardrails}\label{last-mile-guardrails}}

Immediately before submission, the Actuator re-validates every command
against the agent's registered action space and control scope. This
includes commands arriving from a modify verdict of the conflict fabric.
The last mile is the final enforcement point and it does not trust
upstream roles, the fabric included. An out-of-bounds command is refused
with an explicit rejection reason and raises a conformance fault.

This is the third place the same envelope is checked, after admission
and after mediation, and the repetition is intentional. Each check has a
different failure mode available to it, and the last is the only one
positioned after every transformation the intent has undergone.

\hypertarget{declared-failure-behavior-3}{%
\subsection{Declared Failure Behavior}\label{declared-failure-behavior-3}}

How persistent an Actuator is in the face of failure is declared rather
than fixed, because the appropriate persistence differs between a rapid
radio adjustment and a transport reconfiguration. What does not differ
is the set of invariants the declaration operates within: no blind retry
from unknown state, no rollback into unknown state, no automatic retry
of a non-idempotent command whose application is uncertain, and no
acknowledged-unverified outcome reported as success. The policy varies;
the invariants do not.

The failure policy of an Actuator declares the fields of Table 7-1.
Idempotency determines which retry modes are legal at all. The retry
policy carries a maximum attempt count, a window and a backoff, together
with a precondition that is either read-back required or direct, where
direct is legal only for an idempotent command. Unresolved behavior is
either read-back then decide or hold and escalate; blind retry is not an
available value. Unverified behavior is retry with read-back, hold and
escalate, or roll back subject to the read-back precondition. Terminal
failure is either hold or roll back, with an escalation target named.
Rate limits bound actions per window per target.

\begin{wptable}{7-1}

\begin{longtable}[]{@{}
  >{\raggedright\arraybackslash}p{(\columnwidth - 6\tabcolsep) * \real{0.2875}}
  >{\raggedright\arraybackslash}p{(\columnwidth - 6\tabcolsep) * \real{0.1250}}
  >{\raggedright\arraybackslash}p{(\columnwidth - 6\tabcolsep) * \real{0.1000}}
  >{\raggedright\arraybackslash}p{(\columnwidth - 6\tabcolsep) * \real{0.4874}}@{}}
\caption{Declared failure behavior of the Actuator}\tabularnewline
\toprule\noalign{}
\begin{minipage}[b]{\linewidth}\raggedright
\textbf{Field in failure\_\allowbreak{}policy.\allowbreak{}ACT}
\end{minipage} & \begin{minipage}[b]{\linewidth}\raggedright
\textbf{Type}
\end{minipage} & \begin{minipage}[b]{\linewidth}\raggedright
\textbf{Req.}
\end{minipage} & \begin{minipage}[b]{\linewidth}\raggedright
\textbf{Definition}
\end{minipage} \\
\midrule\noalign{}
\endfirsthead
\toprule\noalign{}
\begin{minipage}[b]{\linewidth}\raggedright
\textbf{Field in failure\_\allowbreak{}policy.\allowbreak{}ACT}
\end{minipage} & \begin{minipage}[b]{\linewidth}\raggedright
\textbf{Type}
\end{minipage} & \begin{minipage}[b]{\linewidth}\raggedright
\textbf{Req.}
\end{minipage} & \begin{minipage}[b]{\linewidth}\raggedright
\textbf{Definition}
\end{minipage} \\
\midrule\noalign{}
\endhead
\bottomrule\noalign{}
\endlastfoot
\textbf{idempotency} & enum & M & idempotent or
non\_\allowbreak{}idempotent, per action type. Determines which retry
modes are legal. \\
\textbf{retry\_\allowbreak{}policy} & object & M & Maximum attempts,
window and backoff, with retry\_\allowbreak{}precondition of
readback\_\allowbreak{}required or direct. direct is legal only for an
idempotent action type. \\
\textbf{unresolved\_\allowbreak{}behavior} & enum & M &
readback\_\allowbreak{}then\_\allowbreak{}decide or
hold\_\allowbreak{}and\_\allowbreak{}escalate.\allowbreak{} Blind retry
is not an available value. \\
\textbf{unverified\_\allowbreak{}behavior} & enum & M &
retry\_\allowbreak{}readback,
hold\_\allowbreak{}and\_\allowbreak{}escalate or rollback, the last
subject to the read-back precondition. \\
\textbf{terminal\_\allowbreak{}failure} & object & M & hold or rollback,
with the escalation target named. \\
\textbf{rate\_\allowbreak{}limits} & object & M & Maximum actions per
window per target. Held and applied by the actuation path. \\
\end{longtable}

\end{wptable}

An aggressive agent declares more attempts and a tighter backoff; a
cautious one declares read-back-gated single attempts. Both are
conformant. Neither can declare its way past an invariant, because the
invariants are enforced by the runtime rather than expressed in the
policy.

\hypertarget{execution-feedback}{%
\subsection{Execution Feedback}\label{execution-feedback}}

Every actuation outcome is published as typed feedback carrying the
intent correlation identifier, the outcome, the configuration applied
where it is known, the read-back evidence, timing and a reason code.

Feedback is emitted by the actuation runtime at transaction termination.
It is not contributed code and not optional. A transaction that
terminates without published feedback is itself a conformance fault, and
the platform detects it because the intent correlation identifier
remains unclosed past the declared actuation timeout. Silence is
therefore a detectable state rather than an absence, which matters
because the failure mode this chapter exists to close is precisely one
that looks like nothing happening.

Routing is dual and is applied by the bus rather than chosen by the
agent. Feedback reaches the originating Decision Engine, where it is
evidence for re-decision within the budget of Chapter 6, and the domain
management function's lifecycle and evaluation path, where it feeds
assurance reporting. A route to the Decision Engine alone would leave
the requester blind to repeated unverified actuation, which is the exact
condition under which an agent appears healthy while achieving nothing
physically.

\hypertarget{the-action-contract}{%
\subsection{The Action Contract}\label{the-action-contract}}

The action contract of an Actuator declares the fields of Table 7-2. It
names the controller binding and its pinned interface version range, the
action space expressed over managed objects and their attributes, the
verification method for each action type, the idempotency of each action
type, the control scope within which the capability may act, and the
staging and rate ceilings that apply. An action type whose verification
method is unverifiable carries the restrictions of Section 7.4.1.

\begin{wptable}{7-2}

\begin{longtable}[]{@{}
  >{\raggedright\arraybackslash}p{(\columnwidth - 6\tabcolsep) * \real{0.2875}}
  >{\raggedright\arraybackslash}p{(\columnwidth - 6\tabcolsep) * \real{0.1250}}
  >{\raggedright\arraybackslash}p{(\columnwidth - 6\tabcolsep) * \real{0.1499}}
  >{\raggedright\arraybackslash}p{(\columnwidth - 6\tabcolsep) * \real{0.4375}}@{}}
\caption{Fields of the action contract}\tabularnewline
\toprule\noalign{}
\begin{minipage}[b]{\linewidth}\raggedright
\textbf{Field in action\_\allowbreak{}contract}
\end{minipage} & \begin{minipage}[b]{\linewidth}\raggedright
\textbf{Type}
\end{minipage} & \begin{minipage}[b]{\linewidth}\raggedright
\textbf{Req.}
\end{minipage} & \begin{minipage}[b]{\linewidth}\raggedright
\textbf{Definition}
\end{minipage} \\
\midrule\noalign{}
\endfirsthead
\toprule\noalign{}
\begin{minipage}[b]{\linewidth}\raggedright
\textbf{Field in action\_\allowbreak{}contract}
\end{minipage} & \begin{minipage}[b]{\linewidth}\raggedright
\textbf{Type}
\end{minipage} & \begin{minipage}[b]{\linewidth}\raggedright
\textbf{Req.}
\end{minipage} & \begin{minipage}[b]{\linewidth}\raggedright
\textbf{Definition}
\end{minipage} \\
\midrule\noalign{}
\endhead
\bottomrule\noalign{}
\endlastfoot
\textbf{contract\_\allowbreak{}id} & string & M & Identifier of the
action contract, with its version. \\
\textbf{controller\_\allowbreak{}binding} & object & M & Controller the
capability actuates through, with the interface and its version. \\
\textbf{compatibility} & object & M & Pinned interface version range. An
observed version outside it is a binding fault and blocks actuation. \\
\textbf{action\_\allowbreak{}space} & object & M & Managed objects and
attributes the capability may change, with permitted ranges, expressed
over managed objects rather than controller operations. \\
\textbf{action\_\allowbreak{}types{[}{]}} & array & M & Per action type:
verification method (ms\_\allowbreak{}observation,
controller\_\allowbreak{}state\_\allowbreak{}query or unverifiable) and
idempotency. \\
\textbf{control\_\allowbreak{}scope} & object & M & Scope within which
the capability may act. Re-validated at the last mile. \\
\textbf{staging} & object & C: declared by action type & Staging
sequence and preconditions for actions applied in steps. \\
\textbf{rate\_\allowbreak{}ceiling} & object & M & Ceiling on actions
per window per target, no higher than the catalog entry permits for the
action type. \\
\end{longtable}

\end{wptable}

Table A-1 gives the descriptor location at which each field of Tables
7-1 and 7-2 is declared.

\hypertarget{the-execution-envelope}{%
\subsection{The Execution Envelope}\label{the-execution-envelope}}

Every terminated transaction publishes the envelope of Table 7-3.

\begin{wptable}{7-3}

\begin{longtable}[]{@{}
  >{\raggedright\arraybackslash}p{(\columnwidth - 6\tabcolsep) * \real{0.2875}}
  >{\raggedright\arraybackslash}p{(\columnwidth - 6\tabcolsep) * \real{0.1250}}
  >{\raggedright\arraybackslash}p{(\columnwidth - 6\tabcolsep) * \real{0.1499}}
  >{\raggedright\arraybackslash}p{(\columnwidth - 6\tabcolsep) * \real{0.4375}}@{}}
\caption{Fields of the execution envelope}\tabularnewline
\toprule\noalign{}
\begin{minipage}[b]{\linewidth}\raggedright
\textbf{Field in envelope}
\end{minipage} & \begin{minipage}[b]{\linewidth}\raggedright
\textbf{Type}
\end{minipage} & \begin{minipage}[b]{\linewidth}\raggedright
\textbf{Req.}
\end{minipage} & \begin{minipage}[b]{\linewidth}\raggedright
\textbf{Definition}
\end{minipage} \\
\midrule\noalign{}
\endfirsthead
\toprule\noalign{}
\begin{minipage}[b]{\linewidth}\raggedright
\textbf{Field in envelope}
\end{minipage} & \begin{minipage}[b]{\linewidth}\raggedright
\textbf{Type}
\end{minipage} & \begin{minipage}[b]{\linewidth}\raggedright
\textbf{Req.}
\end{minipage} & \begin{minipage}[b]{\linewidth}\raggedright
\textbf{Definition}
\end{minipage} \\
\midrule\noalign{}
\endhead
\bottomrule\noalign{}
\endlastfoot
\textbf{transaction\_\allowbreak{}id} & uuid & M & Unique identifier of
the actuation transaction, assigned by the ACT runtime. \\
\textbf{intent\_\allowbreak{}id} & uuid & M & Correlation identifier of
the intent actuated. \\
\textbf{agent\_\allowbreak{}instance\_\allowbreak{}id} & string & M &
Instance that performed the actuation. \\
\textbf{action\_\allowbreak{}type} & string & M & Action type executed,
as declared in the action contract. \\
\textbf{target} & object & M & Managed objects acted upon. \\
\textbf{outcome} & enum & M & acknowledged\_\allowbreak{}verified,
acknowledged\_\allowbreak{}unverified, rejected or unresolved. Written
by the ACT runtime, never by contributed logic. \\
\textbf{applied\_\allowbreak{}configuration} & object & C: where known &
Configuration as applied. Absent where the outcome is unresolved. \\
\textbf{readback\_\allowbreak{}evidence} & object & C: outcome is
verified or unverified & Observation references and read state that
established, or failed to establish, verification. \\
\textbf{verification\_\allowbreak{}method} & enum & M & Method declared
for the action type, including unverifiable. \\
\textbf{attempt\_\allowbreak{}index} & integer & M & Position within the
declared retry budget, held by the runtime. \\
\textbf{reason\_\allowbreak{}code} & string & M & Structured reason for
the outcome. \\
\textbf{submitted\_\allowbreak{}time} & timestamp & M & Time the command
was submitted to the controller. \\
\textbf{terminated\_\allowbreak{}time} & timestamp & M & Time the
transaction reached its outcome. \\
\end{longtable}

\end{wptable}

The outcome field is written by the actuation runtime and never by
contributed logic, on the same principle that governs the verdict field
of the intent envelope: the party whose conformance is being assessed
does not author the assessment. Read-back evidence carries the
observation references that established the verified state, so that a
claim of verification is traceable to the measurements that support it
rather than asserted.

\hypertarget{requirements-4}{%
\subsection{Requirements}\label{requirements-4}}

\begin{requirement}{Req-T23-ACT-01}{MUST}{Actuation is a role, not a privilege tier}
\reqfield{Description / rationale}{%
The ACT shall execute for normative and extension agents through the
same controller path. Containment shall be established by declared
control scope, action-space bounds and conflict mediation, and shall
never be established by a separate or restricted execution route.}
\reqfield{Dependency}{%
Req-T23-FW-02}
\reqfield{Traceability (Forward)}{%
\begin{itemize}
\tightlist
\item
  Single actuation path verified at admission.
\item
  No conformance-dependent routing in the runtime.
\end{itemize}}
\reqfield{UNITY-6G components}{%
\begin{itemize}
\tightlist
\item
  ACT runtime
\item
  Admission (IDMO/DMO)
\end{itemize}}
\reqfield{Linked KPIs}{%
\begin{itemize}
\tightlist
\item
  KPI\#1: Deliver four distributed management entities (MS, AE, DE and
  ACT) deployable in each heterogeneous domain.
\item
  KPI\#41: Reduction of 70\% of deployment time of AI-based
  resource/energy management algorithms.
\end{itemize}}
\end{requirement}

\begin{requirement}{Req-T23-ACT-02}{MUST}{Controller binding and version pinning}
\reqfield{Description / rationale}{%
The ACT shall declare its controller binding and the controller
interface version in its interface profile, with a pinned compatibility
range. An observed interface version outside that range shall be a
binding fault that blocks actuation, and shall not be answered with a
best-effort call.}
\reqfield{Dependency}{%
Req-T23-ACT-01}
\reqfield{Traceability (Forward)}{%
\begin{itemize}
\tightlist
\item
  Checked at activation and on controller version change.
\end{itemize}}
\reqfield{UNITY-6G components}{%
\begin{itemize}
\tightlist
\item
  ACT runtime
\item
  Domain controller
\end{itemize}}
\reqfield{Linked KPIs}{%
\begin{itemize}
\tightlist
\item
  KPI\#5: Real-time reconfiguration of cellular network domain in the
  range smaller than 10 ms from taking the decision.
\item
  KPI\#36: Guarantee 99.99999\% service availability and reliability
  based on effective monitoring and LCM.
\end{itemize}}
\end{requirement}

\begin{requirement}{Req-T23-ACT-03}{MUST}{Acknowledgement is never proof of change}
\reqfield{Description / rationale}{%
Every actuation transaction shall terminate in exactly one of
acknowledged\_\allowbreak{}verified,
acknowledged\_\allowbreak{}unverified, rejected or unresolved. Only
acknowledged\_\allowbreak{}verified shall be reported as success. An
acknowledged\_\allowbreak{}unverified outcome shall never be counted
toward an objective as achieved.}
\reqfield{Dependency}{%
Req-T23-ACT-02}
\reqfield{Traceability (Forward)}{%
\begin{itemize}
\tightlist
\item
  Outcome written by the ACT runtime at transaction termination.
\end{itemize}}
\reqfield{UNITY-6G components}{%
\begin{itemize}
\tightlist
\item
  ACT runtime
\end{itemize}}
\reqfield{Linked KPIs}{%
\begin{itemize}
\tightlist
\item
  KPI\#36: Guarantee 99.99999\% service availability and reliability
  based on effective monitoring and LCM.
\item
  KPI\#5: Real-time reconfiguration of cellular network domain in the
  range smaller than 10 ms from taking the decision.
\end{itemize}}
\end{requirement}

\begin{requirement}{Req-T23-ACT-04}{MUST}{Independent read-back and declared unverifiability}
\reqfield{Description / rationale}{%
Verification shall read network or device state through a path
independent of the controller acknowledgement. An action type whose
state a controller cannot expose shall be declared unverifiable, shall
be admitted only with individually enumerated targets, the catalog rate
ceiling for unverifiable execution, per-execution reporting to the DMO
and elevated evaluation scrutiny, and shall never terminate as
acknowledged\_\allowbreak{}verified.\allowbreak{}}
\reqfield{Dependency}{%
Req-T23-ACT-03}
\reqfield{Traceability (Forward)}{%
\begin{itemize}
\tightlist
\item
  Read-back path registered at admission.
\item
  Unverifiable declarations checked against the catalog.
\end{itemize}}
\reqfield{UNITY-6G components}{%
\begin{itemize}
\tightlist
\item
  ACT runtime
\item
  MS runtime
\item
  Admission (IDMO/DMO)
\end{itemize}}
\reqfield{Linked KPIs}{%
\begin{itemize}
\tightlist
\item
  KPI\#36: Guarantee 99.99999\% service availability and reliability
  based on effective monitoring and LCM.
\item
  KPI\#37: Guarantee 99.99999\% service continuity through AI-driven
  proactive management of cloud-native services.
\end{itemize}}
\end{requirement}

\begin{requirement}{Req-T23-ACT-05}{MUST}{Retry derived from idempotency and read-back}
\reqfield{Description / rationale}{%
The platform shall permit retry after a rejected outcome under the
declared policy; shall require read-back before retry after an
unresolved outcome, retrying only where read-back shows the command did
not apply; shall permit retry after an
acknowledged\_\allowbreak{}unverified outcome only for an idempotent
action type; and shall never automatically retry a non-idempotent action
type whose application is uncertain.}
\reqfield{Dependency}{%
Req-T23-ACT-03}
\reqfield{Traceability (Forward)}{%
\begin{itemize}
\tightlist
\item
  Matrix enforced by the actuation path, not by contributed code.
\end{itemize}}
\reqfield{UNITY-6G components}{%
\begin{itemize}
\tightlist
\item
  ACT runtime
\end{itemize}}
\reqfield{Linked KPIs}{%
\begin{itemize}
\tightlist
\item
  KPI\#36: Guarantee 99.99999\% service availability and reliability
  based on effective monitoring and LCM.
\item
  KPI\#31: Reduce service recovery time below 180 s.
\end{itemize}}
\end{requirement}

\begin{requirement}{Req-T23-ACT-06}{MUST}{Rollback only from known state}
\reqfield{Description / rationale}{%
Rollback shall be permitted only from a state established by read-back.
Rollback from an unresolved outcome shall be prohibited until read-back
resolves actual state, and where read-back cannot resolve it the ACT
shall hold and escalate. Hold and escalate shall be the default where
the declaration is absent or ambiguous.}
\reqfield{Dependency}{%
Req-T23-ACT-05}
\reqfield{Traceability (Forward)}{%
\begin{itemize}
\tightlist
\item
  Precondition enforced by the ACT runtime before any rollback
  actuation.
\end{itemize}}
\reqfield{UNITY-6G components}{%
\begin{itemize}
\tightlist
\item
  ACT runtime
\end{itemize}}
\reqfield{Linked KPIs}{%
\begin{itemize}
\tightlist
\item
  KPI\#36: Guarantee 99.99999\% service availability and reliability
  based on effective monitoring and LCM.
\item
  KPI\#31: Reduce service recovery time below 180 s.
\end{itemize}}
\end{requirement}

\begin{requirement}{Req-T23-ACT-07}{MUST}{Last-mile re-validation}
\reqfield{Description / rationale}{%
Immediately before submission the ACT shall re-validate every command
against the agent\textquotesingle s registered action space and control
scope, including commands arriving from a modify verdict of the conflict
fabric. An out-of-bounds command shall be refused with an explicit
rejection reason and shall raise a conformance fault.}
\reqfield{Dependency}{%
Req-T23-DE-04}
\reqfield{Traceability (Forward)}{%
\begin{itemize}
\tightlist
\item
  Performed by the ACT runtime after mediation and before controller
  submission.
\end{itemize}}
\reqfield{UNITY-6G components}{%
\begin{itemize}
\tightlist
\item
  ACT runtime
\item
  Conflict fabric
\end{itemize}}
\reqfield{Linked KPIs}{%
\begin{itemize}
\tightlist
\item
  KPI\#21: Reduce the performance penalty due to conflicts by 20\% by
  means of DLT, hierarchical algorithms and policy-based conflict
  resolution schemes.
\item
  KPI\#4: Real-time reconfiguration of IEEE 802.11 network domain in
  range smaller than 10 ms from taking the decision.
\end{itemize}}
\end{requirement}

\begin{requirement}{Req-T23-ACT-08}{MUST}{Runtime-held retry counters and rate limits}
\reqfield{Description / rationale}{%
Retry counters, backoff and rate limits shall be held and applied by the
actuation path rather than by contributed code. Attempts exceeding a
declared limit shall be suppressed and a conformance fault recorded
against the agent.}
\reqfield{Dependency}{%
Req-T23-ACT-05}
\reqfield{Traceability (Forward)}{%
\begin{itemize}
\tightlist
\item
  Counters held by the ACT runtime.
\item
  Faults raised to the evaluation function.
\end{itemize}}
\reqfield{UNITY-6G components}{%
\begin{itemize}
\tightlist
\item
  ACT runtime
\item
  Evaluation function
\end{itemize}}
\reqfield{Linked KPIs}{%
\begin{itemize}
\tightlist
\item
  KPI\#4: Real-time reconfiguration of IEEE 802.11 network domain in
  range smaller than 10 ms from taking the decision.
\item
  KPI\#36: Guarantee 99.99999\% service availability and reliability
  based on effective monitoring and LCM.
\end{itemize}}
\end{requirement}

\begin{requirement}{Req-T23-ACT-09}{MUST}{Typed feedback at transaction termination}
\reqfield{Description / rationale}{%
Every terminated transaction shall publish typed feedback carrying the
intent correlation identifier, the outcome, the applied configuration
where known, read-back evidence, timing and a reason code. Feedback
shall be emitted by the actuation runtime and shall not be optional
contributed code.}
\reqfield{Dependency}{%
Req-T23-ACT-03}
\reqfield{Traceability (Forward)}{%
\begin{itemize}
\tightlist
\item
  Emitted by the ACT runtime at termination.
\end{itemize}}
\reqfield{UNITY-6G components}{%
\begin{itemize}
\tightlist
\item
  ACT runtime
\end{itemize}}
\reqfield{Linked KPIs}{%
\begin{itemize}
\tightlist
\item
  KPI\#31: Reduce service recovery time below 180 s.
\item
  KPI\#36: Guarantee 99.99999\% service availability and reliability
  based on effective monitoring and LCM.
\item
  KPI\#5: Real-time reconfiguration of cellular network domain in the
  range smaller than 10 ms from taking the decision.
\end{itemize}}
\end{requirement}

\begin{requirement}{Req-T23-ACT-10}{MUST}{Unclosed transactions are conformance faults}
\reqfield{Description / rationale}{%
A transaction that terminates without published feedback shall be a
conformance fault, detected by the platform through an intent
correlation identifier that remains unclosed past the declared actuation
timeout.}
\reqfield{Dependency}{%
Req-T23-ACT-09}
\reqfield{Traceability (Forward)}{%
\begin{itemize}
\tightlist
\item
  Detected by the platform from unclosed correlation identifiers.
\end{itemize}}
\reqfield{UNITY-6G components}{%
\begin{itemize}
\tightlist
\item
  ACT runtime
\item
  DMO
\end{itemize}}
\reqfield{Linked KPIs}{%
\begin{itemize}
\tightlist
\item
  KPI\#31: Reduce service recovery time below 180 s.
\item
  KPI\#36: Guarantee 99.99999\% service availability and reliability
  based on effective monitoring and LCM.
\end{itemize}}
\end{requirement}

\begin{requirement}{Req-T23-ACT-11}{MUST}{Dual routing applied by the bus}
\reqfield{Description / rationale}{%
Execution feedback shall be routed both to the originating DE as
re-decision evidence and to the DMO lifecycle and evaluation path for
assurance reporting. Routing shall be applied by the bus and shall not
be a choice of the emitting agent.}
\reqfield{Dependency}{%
Req-T23-ACT-09}
\reqfield{Traceability (Forward)}{%
\begin{itemize}
\tightlist
\item
  Routing configured at activation by the DMO.
\end{itemize}}
\reqfield{UNITY-6G components}{%
\begin{itemize}
\tightlist
\item
  SMB
\item
  DMO
\item
  DE runtime
\end{itemize}}
\reqfield{Linked KPIs}{%
\begin{itemize}
\tightlist
\item
  KPI\#31: Reduce service recovery time below 180 s.
\item
  KPI\#26: Increase the ratio of service management and LCM tasks
  resolved by local AE/DE components.
\item
  KPI\#1: Deliver four distributed management entities (MS, AE, DE and
  ACT) deployable in each heterogeneous domain.
\item
  KPI\#41: Reduction of 70\% of deployment time of AI-based
  resource/energy management algorithms.
\end{itemize}}
\end{requirement}

\begin{requirement}{Req-T23-ACT-12}{MUST}{Delegated actuation carries the same obligations}
\reqfield{Description / rationale}{%
Where an ACT capability is declared delegated, the delegate shall be
named in the descriptor and shall itself satisfy the requirements of
this chapter. Delegation shall transfer the execution and shall not
transfer or waive any obligation.}
\reqfield{Dependency}{%
Req-T23-ACT-01}
\reqfield{Traceability (Forward)}{%
\begin{itemize}
\tightlist
\item
  Delegate admitted on its own declaration.
\end{itemize}}
\reqfield{UNITY-6G components}{%
\begin{itemize}
\tightlist
\item
  Tuple descriptor
\item
  Admission (IDMO/DMO)
\end{itemize}}
\reqfield{Linked KPIs}{%
\begin{itemize}
\tightlist
\item
  KPI\#1: Deliver four distributed management entities (MS, AE, DE and
  ACT) deployable in each heterogeneous domain.
\item
  KPI\#41: Reduction of 70\% of deployment time of AI-based
  resource/energy management algorithms.
\end{itemize}}
\end{requirement}

\hypertarget{contribution-to-the-project-kpis-3}{%
\subsection{CONTRIBUTION TO THE PROJECT KPIs}\label{contribution-to-the-project-kpis-3}}

Table 7-4 states what the Actuator contributes to each KPI it bears on,
and Table 7-5 maps the requirements of this chapter onto those KPIs.

\begin{wptable}{7-4}

\begin{longtable}[]{@{}
  >{\raggedright\arraybackslash}p{(\columnwidth - 4\tabcolsep) * \real{0.1499}}
  >{\raggedright\arraybackslash}p{(\columnwidth - 4\tabcolsep) * \real{0.6375}}
  >{\raggedright\arraybackslash}p{(\columnwidth - 4\tabcolsep) * \real{0.2126}}@{}}
\caption{Contribution of the Actuator to the project
KPIs}\tabularnewline
\toprule\noalign{}
\begin{minipage}[b]{\linewidth}\raggedright
\textbf{KPI}
\end{minipage} & \begin{minipage}[b]{\linewidth}\raggedright
\textbf{Contribution of the Actuator}
\end{minipage} & \begin{minipage}[b]{\linewidth}\raggedright
\textbf{Requirements}
\end{minipage} \\
\midrule\noalign{}
\endfirsthead
\toprule\noalign{}
\begin{minipage}[b]{\linewidth}\raggedright
\textbf{KPI}
\end{minipage} & \begin{minipage}[b]{\linewidth}\raggedright
\textbf{Contribution of the Actuator}
\end{minipage} & \begin{minipage}[b]{\linewidth}\raggedright
\textbf{Requirements}
\end{minipage} \\
\midrule\noalign{}
\endhead
\bottomrule\noalign{}
\endlastfoot
\textbf{KPI\#1} & Supplies the ACT of the four-entity tuple, executing
for normative and extension agents through one controller path. &
Req-T23-ACT-01, Req-T23-ACT-11 \\
\textbf{KPI\#4} & Last-mile bound re-validation and runtime-held rate
limiting sit on the sub-10 ms reconfiguration path without a separate
execution route. & Req-T23-ACT-07, Req-T23-ACT-08 \\
\textbf{KPI\#5} & The same verification and feedback contract governs
cellular reconfiguration. & Req-T23-ACT-03, Req-T23-ACT-09 \\
\textbf{KPI\#21} & Conflict-driven penalty is reduced because commands
from a modify verdict are re-validated before submission rather than
trusted. & Req-T23-ACT-07 \\
\textbf{KPI\#31} & Typed feedback at transaction termination, with
unclosed correlation identifiers detected as faults, keeps recovery
within the budget rather than waiting on silence. & Req-T23-ACT-09,
Req-T23-ACT-10 \\
\textbf{KPI\#36} & Verification against independently observed state
makes an actuation outcome evidence rather than an assertion, and the
retry and rollback matrix prevents double application and rollback into
unknown state, the two failure modes that turn an incident into an
outage. & Req-T23-ACT-03, Req-T23-ACT-04, Req-T23-ACT-05,
Req-T23-ACT-06 \\
\textbf{KPI\#37} & Unverifiable actions are declared and bounded rather
than invisible, so proactive action is never credited with an effect it
did not demonstrate. & Req-T23-ACT-04 \\
\textbf{KPI\#41} & Extension agents actuate through the same path at
declared bounds, requiring no bespoke execution integration. &
Req-T23-ACT-11 \\
\end{longtable}

\end{wptable}

\begin{wptable}{7-5}\footnotesize

\begin{longtable}[]{@{}
  >{\raggedright\arraybackslash}p{(\columnwidth - 18\tabcolsep) * \real{0.1813}}
  >{\raggedright\arraybackslash}p{(\columnwidth - 18\tabcolsep) * \real{0.0910}}
  >{\raggedright\arraybackslash}p{(\columnwidth - 18\tabcolsep) * \real{0.0910}}
  >{\raggedright\arraybackslash}p{(\columnwidth - 18\tabcolsep) * \real{0.0910}}
  >{\raggedright\arraybackslash}p{(\columnwidth - 18\tabcolsep) * \real{0.0910}}
  >{\raggedright\arraybackslash}p{(\columnwidth - 18\tabcolsep) * \real{0.0910}}
  >{\raggedright\arraybackslash}p{(\columnwidth - 18\tabcolsep) * \real{0.0910}}
  >{\raggedright\arraybackslash}p{(\columnwidth - 18\tabcolsep) * \real{0.0910}}
  >{\raggedright\arraybackslash}p{(\columnwidth - 18\tabcolsep) * \real{0.0910}}
  >{\raggedright\arraybackslash}p{(\columnwidth - 18\tabcolsep) * \real{0.0911}}@{}}
\caption{Mapping of Actuator requirements onto the project
KPIs}\tabularnewline
\toprule\noalign{}
\begin{minipage}[b]{\linewidth}\raggedright
\textbf{KPI\textbackslash{}}

\textbf{Requirement}
\end{minipage} & \begin{minipage}[b]{\linewidth}\raggedright
\textbf{\#1}
\end{minipage} & \begin{minipage}[b]{\linewidth}\raggedright
\textbf{\#4}
\end{minipage} & \begin{minipage}[b]{\linewidth}\raggedright
\textbf{\#5}
\end{minipage} & \begin{minipage}[b]{\linewidth}\raggedright
\textbf{\#21}
\end{minipage} & \begin{minipage}[b]{\linewidth}\raggedright
\textbf{\#26}
\end{minipage} & \begin{minipage}[b]{\linewidth}\raggedright
\textbf{\#31}
\end{minipage} & \begin{minipage}[b]{\linewidth}\raggedright
\textbf{\#36}
\end{minipage} & \begin{minipage}[b]{\linewidth}\raggedright
\textbf{\#37}
\end{minipage} & \begin{minipage}[b]{\linewidth}\raggedright
\textbf{\#41}
\end{minipage} \\
\midrule\noalign{}
\endfirsthead
\toprule\noalign{}
\begin{minipage}[b]{\linewidth}\raggedright
\textbf{KPI\textbackslash{}}

\textbf{Requirement}
\end{minipage} & \begin{minipage}[b]{\linewidth}\raggedright
\textbf{\#1}
\end{minipage} & \begin{minipage}[b]{\linewidth}\raggedright
\textbf{\#4}
\end{minipage} & \begin{minipage}[b]{\linewidth}\raggedright
\textbf{\#5}
\end{minipage} & \begin{minipage}[b]{\linewidth}\raggedright
\textbf{\#21}
\end{minipage} & \begin{minipage}[b]{\linewidth}\raggedright
\textbf{\#26}
\end{minipage} & \begin{minipage}[b]{\linewidth}\raggedright
\textbf{\#31}
\end{minipage} & \begin{minipage}[b]{\linewidth}\raggedright
\textbf{\#36}
\end{minipage} & \begin{minipage}[b]{\linewidth}\raggedright
\textbf{\#37}
\end{minipage} & \begin{minipage}[b]{\linewidth}\raggedright
\textbf{\#41}
\end{minipage} \\
\midrule\noalign{}
\endhead
\bottomrule\noalign{}
\endlastfoot
\textbf{Req-T23-ACT-01} & \kpimark{} & & & & & & & & \kpimark{} \\
\textbf{Req-T23-ACT-02} & & & \kpimark{} & & & & \kpimark{} & & \\
\textbf{Req-T23-ACT-03} & & & \kpimark{} & & & & \kpimark{} & & \\
\textbf{Req-T23-ACT-04} & & & & & & & \kpimark{} & \kpimark{} & \\
\textbf{Req-T23-ACT-05} & & & & & & \kpimark{} & \kpimark{} & & \\
\textbf{Req-T23-ACT-06} & & & & & & \kpimark{} & \kpimark{} & & \\
\textbf{Req-T23-ACT-07} & & \kpimark{} & & \kpimark{} & & & & & \\
\textbf{Req-T23-ACT-08} & & \kpimark{} & & & & & \kpimark{} & & \\
\textbf{Req-T23-ACT-09} & & & \kpimark{} & & & \kpimark{} & \kpimark{} &
& \\
\textbf{Req-T23-ACT-10} & & & & & & \kpimark{} & \kpimark{} & & \\
\textbf{Req-T23-ACT-11} & \kpimark{} & & & & \kpimark{} & \kpimark{} & &
& \kpimark{} \\
\textbf{Req-T23-ACT-12} & \kpimark{} & & & & & & & & \kpimark{} \\
\end{longtable}

\end{wptable}
